\begin{center}
    \textit{\huge{Abstract} }
\end{center}
    
\noindent \textit{The market is constantly searching for professionals skilled in control techniques, making it necessary to introduce mechatronics engineering students to practical control experiences. In this context, the objective of this work emerges: develop a didactic bench for a linear positioner and implement a position control system, developing and validating a mathematical model for it, and comparing simulation and experimental results. The mathematical model was developed in state space, and step response tests and motor torque tests were used to obtain the modeling coefficients. State Feedback was chosen to perform the control task. In this work, an Arduino Mega was used to incorporate the control law, making it possible to vary the average voltage applied to the DC motor via an H-bridge. The position reference considered was of the step type. In both simulation and experiment, the car reached the desired position; however, in the transient regime, the responses become distinct, causing differences from the stipulated design requirements.These discrepancies can be attributed to nonlinear effects present in the real equipment, such as the dead zone and dry friction, as well as uncertainties in the model parameters and the nonlinear nature of some quantities. However, the control system achieved its main objective, which is to guide the car to the desired position.}

\vspace{\onelineskip}
    
\noindent
\textbf{Keywords}: \textit{Linear positioner}; \textit{Didactic bench}; \textit{State space}; \textit{State-feedback control}; \textit{Closed-loop control}.

\newpage